%%=============================================================================
%% CAPITULO 1 -- INTRODUCAO
%%
%% Na introducao, faca a apresentacao do trabalho. Ela devera ocupar uma ou,
%% no maximo, duas paginas. De um panorama da tematica, exponha a questao
%% problema que conduzira a pesquisa proposta, a hipotese e a justificativa.
%%
%% Titulos: nao utilize ponto apos a numeracao das secoes.
%%=============================================================================

\section{Introdução}

Sistemas de Retrieval-Augmented Generation (RAG) respondem perguntas
combinando um modelo de linguagem com uma busca sobre uma base documental, o
que reduz alucinação e permite atualizar o conhecimento do sistema sem
retreinar o modelo \parencite{lewis2020}. Essa vantagem depende da qualidade
da recuperação. Em corpus jurídico-administrativo brasileiro --- leis
federais, resoluções do Conselho Nacional de Justiça e instruções normativas
do Tribunal de Contas da União ---, a busca vetorial pura
erra na recuperação de termos exatos (siglas, números de processo, remissões
normativas) e não resolve perguntas que exigem conectar informação
distribuída entre documentos distintos, como saber se uma norma continua em
vigor diante de um ato que a revoga. O GraphRAG ataca essa segunda limitação
complementando a busca vetorial com um grafo de conhecimento entre documentos
\parencite{edge2024}, mas as implementações consolidadas reconstroem o grafo
inteiro em lote, premissa incompatível com um acervo que cresce
continuamente.

\subsection{Objetivo geral}

Comparar empiricamente, sobre um corpus jurídico-administrativo brasileiro em
português, estratégias de recuperação de informação (RAG denso, RAG
denso-esparso e GraphRAG incremental) e estratégias de geração de resposta
(single-pass e verificada multiagente).

\subsection{Objetivos específicos}

\begin{itemize}
  \item Construir e curar o corpus jurídico-administrativo e desenvolver um harness de avaliação com golden dataset de perguntas e respostas de referência;
  \item Avaliar se o GraphRAG incremental supera RAG denso e RAG denso-esparso na recuperação de evidências e na qualidade das respostas geradas (H1);
  \item Avaliar se a geração verificada multiagente melhora a fidelidade e a rastreabilidade das respostas em relação à geração single-pass (H2).
\end{itemize}

\subsection{Problema}

Diante da limitação da busca vetorial pura para conectar informação
distribuída entre documentos jurídico-administrativos e da incompatibilidade
das implementações consolidadas de GraphRAG com um acervo que cresce
continuamente, qual estratégia de recuperação --- RAG denso, RAG
denso-esparso ou GraphRAG incremental --- e qual estratégia de geração ---
single-pass ou verificada multiagente --- produz respostas mais precisas,
fundamentadas e rastreáveis sobre um corpus jurídico-administrativo
brasileiro?

\subsection{Hipótese}

\textbf{H1}: a condição GraphRAG Incremental (pgvector + Apache AGE) supera
as condições RAG Denso e RAG Denso-Esparso em Context Precision, Context
Recall, Faithfulness e Answer Relevancy, avaliada sobre um golden dataset com
perguntas single-hop e multi-hop, mantendo idêntico o pipeline de ingestão e
embedding entre as três condições.

\textbf{H2}: a geração verificada multiagente (padrão
Researcher-Auditor-Adjudicator), aplicada sobre a estratégia de recuperação
vencedora de H1, melhora a Faithfulness e a taxa de verificação de citação em
relação à geração single-pass sobre a mesma evidência recuperada.

\subsection{Justificativa}

Do ponto de vista arquitetural, nenhum trabalho revisado combina storage
único relacional, atualização incremental por documento e ausência de
sumarização de comunidades: o GraphRAG canônico \parencite{edge2024}
reconstrói o grafo em lote, e as alternativas que atualizam incrementalmente
o fazem sem storage unificado. Do ponto de vista empírico, os benchmarks de
recuperação jurídica em português revisados reúnem jurisprudência e normas em
coleções separadas e avaliam apenas a recuperação, não a resposta gerada. Um
acervo que combine legislação federal e atos normativos de órgãos de mais de
um poder, avaliado de ponta a ponta, da recuperação à geração, carece de
referência dedicada.
